\subsection*{Ekwalizatory}

Niech $\C$ będzie kategorią.

\begin{definicja}
Dla pary strzałek $f, g \colon A \to B \in \C$ mówimy, że $e\colon E\to A$ jest \emph{ekwalizatorem} $f$ i $g$, jeśli:
\begin{enumerate}
\item diagram
\[\begin{tikzcd}
E \arrow[r,"e"] & A \arrow[r,shift left,"f"] \arrow[r,shift right,swap,"g"] & B
\end{tikzcd}\]
jest przemienny, tj.\ $f\circ e = g\circ e$,
\item dla dowolnego $z\colon Z\to A$ takiego, że $f\circ z = g\circ z$, istnieje dokładnie jedna strzałka $u\colon Z\to E$ czyniąca następujący diagram przemiennym:
\[\begin{tikzcd}
E \arrow[r,"e"] & A \arrow[r,shift left,"f"] \arrow[r,shift right,swap,"g"] & B \\
Z \arrow[ur,swap,"z"] \arrow[u,dashed,"\exists !\, u"] & &
\end{tikzcd}\]
\end{enumerate}
\end{definicja}

\begin{przyklad}
W kategorii $\Set$, mając dane dwa przekształcenia $f,g\colon A\to B$, ich ekwalizatorem będzie
\[
\{x\in A \mid f(x)=g(x)\} \hookrightarrow A,
\]
gdzie strzałka jest zanurzeniem (inkluzją).
\end{przyklad}

\paragraph{Dowód.}
Dla każdego $z\colon Z\to A$ takiego, że $f\circ z = g\circ z$, mamy: dla $y\in Z$ zachodzi $f(z(y)) = g(z(y))$, czyli element $x:=z(y)\in A$ spełnia $f(x)=g(x)$. Definiujemy
\[
u\colon Z \to \{x\in A\mid f(x)=g(x)\}, \qquad u(y) := z(y).
\]
Wtedy diagram
\[\begin{tikzcd}
\{x\in A \mid f(x)=g(x)\} \arrow[r,hook] & A \arrow[r,shift left,"f"] \arrow[r,shift right,swap,"g"] & B \\
Z \arrow[ur,swap,"z"] \arrow[u,"u"] & &
\end{tikzcd}\]
jest przemienny. Przekształcenie $u$ jest jednoznaczne. \qed

\begin{stwierdzenie}
Jeśli $e\colon E\to A$ jest ekwalizatorem pary $A\rightrightarrows B$, to $e$ jest mono.
\end{stwierdzenie}

\begin{przyklad}
Rozważmy w $\Set$ podzbiór $U\subseteq A$ oraz diagram
\[\begin{tikzcd}
? \arrow[r,"?"] & A \arrow[r,shift left,"\mathrm{const}_1"] \arrow[r,shift right,swap,"\chi_U"] & 2=\{0,1\}
\end{tikzcd}\]
gdzie
\[
\chi_U(a)=\begin{cases} 1 & a\in U,\\ 0 & \text{w p.p.},\end{cases} \qquad \mathrm{const}_1(a)=1.
\]
Czym jest ekwalizator $\mathrm{const}_1$ i $\chi_U$? Warunek $\chi_U(a) = \mathrm{const}_1(a) = 1$ oznacza dokładnie $a\in U$, więc ekwalizatorem jest inkluzja $U \hookrightarrow A$. W szczególności \emph{każdy} podzbiór $U\subseteq A$ jest ekwalizatorem pewnej pary strzałek o wartościach w $2=\{0,1\}$.
\end{przyklad}

\subsection*{Koekwalizatory}

\begin{definicja}
Dla pary strzałek $f, g \colon A\to B \in \C$ \emph{koekwalizatorem} $f$ i $g$ nazywamy strzałkę $q\colon B\to Q$ taką, że:
\begin{enumerate}
\item $q\circ f = q\circ g$, tj.\ diagram
\[\begin{tikzcd}
A \arrow[r,shift left,"f"] \arrow[r,shift right,swap,"g"] & B \arrow[r,"q"] & Q
\end{tikzcd}\]
jest przemienny,
\item dla każdego $z\colon B\to Z$ takiego, że $z\circ f = z\circ g$, istnieje dokładnie jedna strzałka $u\colon Q\to Z$ taka, że poniższy diagram jest przemienny:
\[\begin{tikzcd}
A \arrow[r,shift left,"f"] \arrow[r,shift right,swap,"g"] & B \arrow[r,"q"] \arrow[dr,swap,"z"] & Q \arrow[d,dashed,"\exists !\, u"] \\
& & Z
\end{tikzcd}\]
\end{enumerate}
\end{definicja}

\begin{stwierdzenie}
Jeśli $q\colon B\to Q$ jest koekwalizatorem pewnej pary strzałek, to $q$ jest epi.
\end{stwierdzenie}

\begin{przyklad}
Niech $R\subseteq X\times X$ będzie relacją binarną na $X$. Rozważmy diagram
\[\begin{tikzcd}
R \arrow[r,shift left,"\pi_1"] \arrow[r,shift right,swap,"\pi_2"] & X
\end{tikzcd}\]
Jeśli $R$ jest relacją równoważności, to koekwalizatorem tej pary rzutowań jest
\[\begin{tikzcd}
R \arrow[r,shift left,"\pi_1"] \arrow[r,shift right,swap,"\pi_2"] & X \arrow[r,"\mathrm{nat}\,R"] & X/R,
\end{tikzcd}\]
gdzie
\[
X/R = \{\,[x]_R \mid x\in X\,\}
\]
(klasy abstrakcji relacji $R$),
\[
\mathrm{nat}\,R\colon X \twoheadrightarrow X/R, \qquad \mathrm{nat}\,R(x) \stackrel{\mathrm{df}}{=} [x]_R.
\]
\end{przyklad}

\subsection*{Pullback}

\begin{definicja}
\emph{Pullbackiem} (cofnięciem) pary strzałek
\[\begin{tikzcd}
& B \arrow[d,"g"] \\
A \arrow[r,"f"] & C
\end{tikzcd}\]
nazywamy obiekt $P$ wraz ze strzałkami $p_1\colon P\to A$, $p_2\colon P\to B$ takimi, że
\begin{enumerate}
\item diagram
\[\begin{tikzcd}
P \arrow[r,"p_2"] \arrow[d,swap,"p_1"] & B \arrow[d,"g"] \\
A \arrow[r,"f"] & C
\end{tikzcd}\]
jest przemienny,
\item dla każdego obiektu $Z$ wraz ze strzałkami $z_1\colon Z\to A$ oraz $z_2\colon Z\to B$ takimi, że $f\circ z_1 = g\circ z_2$, istnieje dokładnie jedna strzałka $u\colon Z\to P$ czyniąca następujący diagram przemiennym:
\[\begin{tikzcd}
Z \arrow[drr,bend left,"z_2"] \arrow[ddr,bend right,swap,"z_1"] \arrow[dr,dashed,"\exists !\, u"] & & \\
& P \arrow[r,"p_2"] \arrow[d,swap,"p_1"] & B \arrow[d,"g"] \\
& A \arrow[r,swap,"f"] & C
\end{tikzcd}\]
\end{enumerate}
\end{definicja}

\begin{uwaga}
Często używać będziemy notacji ,,produktowej'' w kontekście pullbacków. Dokładniej, obiekt pullback strzałek $A\xrightarrow{f} C \xleftarrow{g} B$ oznaczać będziemy przez $A\times_C B$. Morfizmy pullbacka często będziemy też oznaczać przez $\pi_1,\pi_2$:
\[\begin{tikzcd}
A\times_C B \arrow[r,"\pi_2"] \arrow[d,swap,"\pi_1"] & B \arrow[d,"g"] \\
A \arrow[r,swap,"f"] & C
\end{tikzcd}\]
\end{uwaga}

\begin{przyklad}
W $\Set$ pullbackiem pary $A\xrightarrow{f} C \xleftarrow{g} B$ jest
\[
A\times_C B \;=\; \{(a,b)\in A\times B \;:\; f(a) = g(b)\}
\]
wraz z obcięciami rzutowań. Dwa ważne przypadki szczególne:
\begin{enumerate}
\item \emph{Przeciwobraz.} Dla inkluzji $i_U \colon U \hookrightarrow C$ podzbioru $U\subseteq C$ mamy
\[
A \times_C U \;=\; \{(a,u) : f(a) = u\} \;\cong\; f^{-1}(U),
\]
czyli pullback wzdłuż $f$ ,,cofa'' podzbiór $U$ do jego przeciwobrazu (por.\ zadanie~\ref{zad:pullback-przeciwobraz}). Stąd właśnie nazwa \emph{cofnięcie}.
\item \emph{Przecięcie.} Dla inkluzji $U \hookrightarrow X \hookleftarrow V$ dwóch podzbiorów $U, V \subseteq X$ pullbackiem jest $U \cap V$ wraz z inkluzjami.
\end{enumerate}
\end{przyklad}

\begin{stwierdzenie}
Niech $\C$ będzie kategorią z binarnymi produktami i ekwalizatorami. Dla danej pary $A\xrightarrow{f} C \xleftarrow{g} B$:
\begin{enumerate}
\item weźmy produkt binarny $A\times B$ wraz z $\pi_1\colon A\times B\to A$, $\pi_2\colon A\times B\to B$,
\item rozważmy ekwalizator $E \xhookrightarrow{e} A\times B$ pary
\[\begin{tikzcd}
A\times B \arrow[r,shift left,"f\circ \pi_1"] \arrow[r,shift right,swap,"g\circ \pi_2"] & C.
\end{tikzcd}\]
\end{enumerate}
Wtedy obiekt $E$ wraz z $\pi_1\circ e\colon E\to A$ oraz $\pi_2\circ e\colon E\to B$ jest pullbackiem $A\xrightarrow{f} C \xleftarrow{g} B$.
\end{stwierdzenie}

\paragraph{Dowód.}
Widzimy, że $E$ wraz z $\pi_1\circ e$ oraz $\pi_2\circ e$ czyni kwadrat $f\circ(\pi_1\circ e) = g\circ(\pi_2\circ e)$ przemiennym -- wynika to z faktu, że $e$ jest ekwalizatorem $f\circ\pi_1$ oraz $g\circ\pi_2$.

Weźmy obiekt $Z$ wraz z $z_1\colon Z\to A$, $z_2\colon Z\to B$ takie, że $f\circ z_1 = g\circ z_2$. Z własności produktu istnieje jednoznaczna strzałka $\langle z_1,z_2\rangle\colon Z\to A\times B$ z $\pi_1\circ\langle z_1,z_2\rangle = z_1$, $\pi_2\circ\langle z_1,z_2\rangle = z_2$. Widzimy, że
\[
f\circ \pi_1 \circ \langle z_1,z_2\rangle = f\circ z_1 = g\circ z_2 = g\circ \pi_2\circ\langle z_1,z_2\rangle,
\]
więc istnieje dokładnie jedna strzałka $u\colon Z\to E$ taka, że $e\circ u = \langle z_1,z_2\rangle$. Wówczas
\[
(\pi_1\circ e)\circ u = z_1, \qquad (\pi_2\circ e)\circ u = z_2.
\]
Pozostaje jednoznaczność: jeśli $u'\colon Z\to E$ również spełnia $(\pi_1\circ e)\circ u' = z_1$ oraz $(\pi_2\circ e)\circ u' = z_2$, to z jednoznaczności strzałki do produktu $e\circ u' = \langle z_1,z_2\rangle$, a stąd, z jednoznaczności w definicji ekwalizatora, $u' = u$. \qed
